
\documentclass{article}
\usepackage{colortbl}
\usepackage{makecell}
\usepackage{multirow}
\usepackage{supertabular}

\begin{document}

\newcounter{utterance}

\centering \large Interaction Transcript for game `dond', experiment `semi\_en', episode 6 with reflexion{-}Qwen3.5{-}4B.
\vspace{24pt}

{ \footnotesize  \setcounter{utterance}{1}
\setlength{\tabcolsep}{0pt}
\begin{supertabular}{c@{$\;$}|p{.15\linewidth}@{}p{.15\linewidth}p{.15\linewidth}p{.15\linewidth}p{.15\linewidth}p{.15\linewidth}}
    \# & $\;$A & \multicolumn{4}{c}{Game Master} & $\;\:$B\\
    \hline

    \theutterance \stepcounter{utterance}  
    & & \multicolumn{4}{p{0.6\linewidth}}{
        \cellcolor[rgb]{0.9,0.9,0.9}{
            \makecell[{{p{\linewidth}}}]{
                \texttt{\tiny{[P1$\langle$GM]}}
                \texttt{You are playing a negotiation game in which you have to agree on how to divide a set of items among you and another player.} \\
\\ 
\texttt{Rules:} \\
\texttt{(a) You and the other player are given a set of items. Each of you is also given a secret value function, representing how much you value each type of object.} \\
\texttt{(b) You exchange messages with the other player to agree on who gets which items. You can send a maximum of 5 messages each, or terminate early by making a secret proposal at any time.} \\
\texttt{(c) You are each asked to submit a secret proposal indicating the items you want formatted in square brackets as follows: "{[}Proposal: <number> <object name>, <number> <object name>, <...>{]}"} \\
\texttt{(d) If your proposals are complementary, i.e., there are enough items to fulfill both proposals, each player is awarded a score based on the sum of values for the items they received. Otherwise, both of you get zero points.} \\
\texttt{(e) Your goal is to maximize the score you receive.} \\
\\ 
\texttt{Let us start.} \\
\\ 
\texttt{The set of available items is:} \\
\\ 
\texttt{2 eggs, 1 picture, 2 dice} \\
\\ 
\texttt{Your secret value function is:} \\
\\ 
\texttt{egg: 1, picture: 2, dice: 3} \\
\\ 
\texttt{IMPORTANT: Your messages, unless it is the secret proposal, are directly transmitted to the other player, so do not include any response to the rules or text announcing your message. To make a secret proposal, use the indicated format. Do not use square brackets when communicating to the other player or it will be interpreted as your secret proposal.} \\
            }
        }
    }
    & \\ \\

    \theutterance \stepcounter{utterance}  
    & \multicolumn{4}{p{0.6\linewidth}}{
        \cellcolor[rgb]{0.9,0.9,0.9}{
            \makecell[{{p{\linewidth}}}]{
                \texttt{\tiny{[P1$\rangle$GM]}}
                \texttt{I will propose a division that maximizes my value while ensuring the other player also receives a positive value (to avoid the zero{-}point penalty).} \\
\\ 
\texttt{My values:} \\
\texttt{{-} 2 eggs: 2} \\
\texttt{{-} 1 picture: 2} \\
\texttt{{-} 2 dice: 6} \\
\texttt{Total: 10} \\
\\ 
\texttt{If I take the 2 dice (value 6), I leave the other player with 2 eggs and 1 picture.} \\
\texttt{The other player's values (assuming they are rational and want to maximize their own score, though I don't know their specific function, I must assume they value these items positively to accept):} \\
\texttt{{-} 2 eggs: 2} \\
\texttt{{-} 1 picture: 2} \\
\texttt{Total for them: 4.} \\
\\ 
\texttt{This split (Me: 6, Them: 4) is better for me than taking everything (10) because if I take everything, they get 0, and the rule states "If your proposals are complementary... Otherwise, both of you get zero points." Since I cannot know their exact values, I must propose a split where they are likely to accept (i.e., they get something valuable).} \\
\\ 
\texttt{Let's try to give them the items I value least, or a balanced split.} \\
\texttt{Option A: I take 2 dice (6). They take 2 eggs + 1 picture (4).} \\
\texttt{Option B: I take 1 picture + 2 eggs (4). They take 2 dice (6).} \\
\texttt{Option C: I take 2 dice + 1 picture (8). They take 2 eggs (2).} \\
\\ 
\texttt{Option A seems safest as it gives the other player a significant amount (4) without giving me the maximum possible (10), which might trigger a rejection if they value the dice highly. However, since I don't know their values, the safest bet to ensure a non{-}zero outcome is to give them the items I value least, or a fair split.} \\
\\ 
\texttt{Actually, the rule says "If your proposals are complementary... Otherwise... zero points." This implies if I propose "I want X" and they propose "I want Y", and X and Y don't cover all items or overlap in a way that leaves nothing for the other, we get zero. Wait, the rule says "enough items to fulfill both proposals". If I say "I want the 2 dice" and they say "I want the 2 eggs and 1 picture", that is complementary. If I say "I want the 2 dice" and they say "I want the 2 dice", that is not complementary (overlap).} \\
\\ 
\texttt{So I need to propose a set of items that, if accepted by the other player, leaves the rest for them.} \\
\texttt{Let's propose: I want the 2 dice.} \\
\texttt{This leaves 2 eggs and 1 picture for them.} \\
\texttt{If they agree to take the 2 eggs and 1 picture, we succeed} \\
            }
        }
    }
    & & \\ \\

    \theutterance \stepcounter{utterance}  
    & & & \multicolumn{2}{p{0.3\linewidth}}{
        \cellcolor[rgb]{0.9,0.9,0.9}{
            \makecell[{{p{\linewidth}}}]{
                \texttt{\tiny{[GM$|$GM]}}
                \texttt{continue} \\
            }
        }
    }
    & & \\ \\

    \theutterance \stepcounter{utterance}  
    & & & \multicolumn{2}{p{0.3\linewidth}}{
        \cellcolor[rgb]{0.9,0.9,0.9}{
            \makecell[{{p{\linewidth}}}]{
                \texttt{\tiny{[GM$|$GM]}}
                \texttt{I will propose a division that maximizes my value while ensuring the other player also receives a positive value (to avoid the zero{-}point penalty).} \\
\\ 
\texttt{My values:} \\
\texttt{{-} 2 eggs: 2} \\
\texttt{{-} 1 picture: 2} \\
\texttt{{-} 2 dice: 6} \\
\texttt{Total: 10} \\
\\ 
\texttt{If I take the 2 dice (value 6), I leave the other player with 2 eggs and 1 picture.} \\
\texttt{The other player's values (assuming they are rational and want to maximize their own score, though I don't know their specific function, I must assume they value these items positively to accept):} \\
\texttt{{-} 2 eggs: 2} \\
\texttt{{-} 1 picture: 2} \\
\texttt{Total for them: 4.} \\
\\ 
\texttt{This split (Me: 6, Them: 4) is better for me than taking everything (10) because if I take everything, they get 0, and the rule states "If your proposals are complementary... Otherwise, both of you get zero points." Since I cannot know their exact values, I must propose a split where they are likely to accept (i.e., they get something valuable).} \\
\\ 
\texttt{Let's try to give them the items I value least, or a balanced split.} \\
\texttt{Option A: I take 2 dice (6). They take 2 eggs + 1 picture (4).} \\
\texttt{Option B: I take 1 picture + 2 eggs (4). They take 2 dice (6).} \\
\texttt{Option C: I take 2 dice + 1 picture (8). They take 2 eggs (2).} \\
\\ 
\texttt{Option A seems safest as it gives the other player a significant amount (4) without giving me the maximum possible (10), which might trigger a rejection if they value the dice highly. However, since I don't know their values, the safest bet to ensure a non{-}zero outcome is to give them the items I value least, or a fair split.} \\
\\ 
\texttt{Actually, the rule says "If your proposals are complementary... Otherwise... zero points." This implies if I propose "I want X" and they propose "I want Y", and X and Y don't cover all items or overlap in a way that leaves nothing for the other, we get zero. Wait, the rule says "enough items to fulfill both proposals". If I say "I want the 2 dice" and they say "I want the 2 eggs and 1 picture", that is complementary. If I say "I want the 2 dice" and they say "I want the 2 dice", that is not complementary (overlap).} \\
\\ 
\texttt{So I need to propose a set of items that, if accepted by the other player, leaves the rest for them.} \\
\texttt{Let's propose: I want the 2 dice.} \\
\texttt{This leaves 2 eggs and 1 picture for them.} \\
\texttt{If they agree to take the 2 eggs and 1 picture, we succeed} \\
            }
        }
    }
    & & \\ \\

    \theutterance \stepcounter{utterance}  
    & & \multicolumn{4}{p{0.6\linewidth}}{
        \cellcolor[rgb]{0.9,0.9,0.9}{
            \makecell[{{p{\linewidth}}}]{
                \texttt{\tiny{[GM$\rangle$P2]}}
                \texttt{You are playing a negotiation game in which you have to agree on how to divide a set of items among you and another player.} \\
\\ 
\texttt{Rules:} \\
\texttt{(a) You and the other player are given a set of items. Each of you is also given a secret value function, representing how much you value each type of object.} \\
\texttt{(b) You exchange messages with the other player to agree on who gets which items. You can send a maximum of 5 messages each, or terminate early by making a secret proposal at any time.} \\
\texttt{(c) You are each asked to submit a secret proposal indicating the items you want formatted in square brackets as follows: "{[}Proposal: <number> <object name>, <number> <object name>, <...>{]}"} \\
\texttt{(d) If your proposals are complementary, i.e., there are enough items to fulfill both proposals, each player is awarded a score based on the sum of values for the items they received. Otherwise, both of you get zero points.} \\
\texttt{(e) Your goal is to maximize the score you receive.} \\
\\ 
\texttt{Let us start.} \\
\\ 
\texttt{The set of available items is:} \\
\\ 
\texttt{2 eggs, 1 picture, 2 dice} \\
\\ 
\texttt{Your secret value function is:} \\
\\ 
\texttt{egg: 0, picture: 8, dice: 1} \\
\\ 
\texttt{IMPORTANT: Your messages, unless it is the secret proposal, are directly transmitted to the other player, so do not include any response to the rules or text announcing your message. To make a secret proposal, use the indicated format. Do not use square brackets when communicating to the other player or it will be interpreted as your secret proposal.} \\
\\ 
\\ 
\texttt{I will propose a division that maximizes my value while ensuring the other player also receives a positive value (to avoid the zero{-}point penalty).} \\
\\ 
\texttt{My values:} \\
\texttt{{-} 2 eggs: 2} \\
\texttt{{-} 1 picture: 2} \\
\texttt{{-} 2 dice: 6} \\
\texttt{Total: 10} \\
\\ 
\texttt{If I take the 2 dice (value 6), I leave the other player with 2 eggs and 1 picture.} \\
\texttt{The other player's values (assuming they are rational and want to maximize their own score, though I don't know their specific function, I must assume they value these items positively to accept):} \\
\texttt{{-} 2 eggs: 2} \\
\texttt{{-} 1 picture: 2} \\
\texttt{Total for them: 4.} \\
\\ 
\texttt{This split (Me: 6, Them: 4) is better for me than taking everything (10) because if I take everything, they get 0, and the rule states "If your proposals are complementary... Otherwise, both of you get zero points." Since I cannot know their exact values, I must propose a split where they are likely to accept (i.e., they get something valuable).} \\
\\ 
\texttt{Let's try to give them the items I value least, or a balanced split.} \\
\texttt{Option A: I take 2 dice (6). They take 2 eggs + 1 picture (4).} \\
\texttt{Option B: I take 1 picture + 2 eggs (4). They take 2 dice (6).} \\
\texttt{Option C: I take 2 dice + 1 picture (8). They take 2 eggs (2).} \\
\\ 
\texttt{Option A seems safest as it gives the other player a significant amount (4) without giving me the maximum possible (10), which might trigger a rejection if they value the dice highly. However, since I don't know their values, the safest bet to ensure a non{-}zero outcome is to give them the items I value least, or a fair split.} \\
\\ 
\texttt{Actually, the rule says "If your proposals are complementary... Otherwise... zero points." This implies if I propose "I want X" and they propose "I want Y", and X and Y don't cover all items or overlap in a way that leaves nothing for the other, we get zero. Wait, the rule says "enough items to fulfill both proposals". If I say "I want the 2 dice" and they say "I want the 2 eggs and 1 picture", that is complementary. If I say "I want the 2 dice" and they say "I want the 2 dice", that is not complementary (overlap).} \\
\\ 
\texttt{So I need to propose a set of items that, if accepted by the other player, leaves the rest for them.} \\
\texttt{Let's propose: I want the 2 dice.} \\
\texttt{This leaves 2 eggs and 1 picture for them.} \\
\texttt{If they agree to take the 2 eggs and 1 picture, we succeed} \\
            }
        }
    }
    & \\ \\

    \theutterance \stepcounter{utterance}  
    & & & \multicolumn{4}{p{0.6\linewidth}}{
        \cellcolor[rgb]{0.9,0.9,0.9}{
            \makecell[{{p{\linewidth}}}]{
                \texttt{\tiny{[GM$\langle$P2]}}
                \texttt{{[}Proposal: 2 dice{]}} \\
            }
        }
    }
     \\ \\

    \theutterance \stepcounter{utterance}  
    & & & \multicolumn{2}{p{0.3\linewidth}}{
        \cellcolor[rgb]{0.9,0.9,0.9}{
            \makecell[{{p{\linewidth}}}]{
                \texttt{\tiny{[GM$|$GM]}}
                \texttt{proposal} \\
            }
        }
    }
    & & \\ \\

    \theutterance \stepcounter{utterance}  
    & & & \multicolumn{2}{p{0.3\linewidth}}{
        \cellcolor[rgb]{0.9,0.9,0.9}{
            \makecell[{{p{\linewidth}}}]{
                \texttt{\tiny{[GM$|$GM]}}
                [0, 0, 2]
            }
        }
    }
    & & \\ \\

    \theutterance \stepcounter{utterance}  
    & & \multicolumn{4}{p{0.6\linewidth}}{
        \cellcolor[rgb]{0.9,0.9,0.9}{
            \makecell[{{p{\linewidth}}}]{
                \texttt{\tiny{[P1$\langle$GM]}}
                \texttt{The other player has submitted a secret proposal. It is now your turn to submit a proposal of your own.} \\
            }
        }
    }
    & \\ \\

    \theutterance \stepcounter{utterance}  
    & \multicolumn{4}{p{0.6\linewidth}}{
        \cellcolor[rgb]{0.9,0.9,0.9}{
            \makecell[{{p{\linewidth}}}]{
                \texttt{\tiny{[P1$\rangle$GM]}}
                \texttt{{[}Proposal: 2 dice{]}} \\
            }
        }
    }
    & & \\ \\

    \theutterance \stepcounter{utterance}  
    & & & \multicolumn{2}{p{0.3\linewidth}}{
        \cellcolor[rgb]{0.9,0.9,0.9}{
            \makecell[{{p{\linewidth}}}]{
                \texttt{\tiny{[GM$|$GM]}}
                \texttt{proposal} \\
            }
        }
    }
    & & \\ \\

    \theutterance \stepcounter{utterance}  
    & & & \multicolumn{2}{p{0.3\linewidth}}{
        \cellcolor[rgb]{0.9,0.9,0.9}{
            \makecell[{{p{\linewidth}}}]{
                \texttt{\tiny{[GM$|$GM]}}
                [0, 0, 2]
            }
        }
    }
    & & \\ \\

    \theutterance \stepcounter{utterance}  
    & & & \multicolumn{2}{p{0.3\linewidth}}{
        \cellcolor[rgb]{0.9,0.9,0.9}{
            \makecell[{{p{\linewidth}}}]{
                \texttt{\tiny{[GM$|$GM]}}
                [[0, 0, 2], [0, 0, 2]]
            }
        }
    }
    & & \\ \\

    \theutterance \stepcounter{utterance}  
    & \multicolumn{4}{p{0.6\linewidth}}{
        \cellcolor[rgb]{0.95,0.95,0.95}{
            \makecell[{{p{\linewidth}}}]{
                \texttt{\tiny{[P2$\rightarrow$P1]}}
                \texttt{Post{-}episode Reflexion (not sent to GM):} \\
\texttt{New plan: I failed because I only proposed a single item (the 2 dice) without explicitly stating what I wanted the other player to take, or without ensuring my proposal was clearly complementary to a likely counter{-}proposal. The rules state that success depends on proposals being "complementary" (covering all items without overlap). By only saying "I want the 2 dice," I left the other player with 2 eggs and 1 picture, but I did not formally propose that specific remainder for them, nor did I account for the possibility that they might want the dice. To succeed, I must propose a complete partition of the items in my message. Specifically, I should propose taking the 2 dice and explicitly state that the other player gets the 2 eggs and 1 picture (e.g., "{[}Proposal: 2 dice, 2 eggs, 1 picture{]}"), or propose a split where I take the items I value least (the eggs and picture) and they take the dice, ensuring the sum of our values is maximized and the sets are perfectly disjoint. I will calculate the optimal split based on my known values (Dice=6, Picture=2, Eggs=2) and propose the exact division of all items to guarantee complementarity.} \\
            }
        }
    }
    & & \\ \\

    \theutterance \stepcounter{utterance}  
    & & & \multicolumn{4}{p{0.6\linewidth}}{
        \cellcolor[rgb]{0.95,0.95,0.95}{
            \makecell[{{p{\linewidth}}}]{
                \texttt{\tiny{[P1$\rightarrow$P2]}}
                \texttt{Post{-}episode Reflexion (not sent to GM):} \\
\texttt{New plan: I failed because I submitted a secret proposal immediately without first communicating with the other player to gauge their preferences or reach a mutual agreement through dialogue. The rules state that proposals are only successful if they are "complementary" and accepted by both parties; a unilateral proposal without prior negotiation likely resulted in a mismatch or rejection. Next time, I will start by sending a message to the other player proposing a specific division (e.g., "I propose you take the 2 eggs and 1 picture, and I take the 2 dice") to see if they agree. Only if we reach a consensus through conversation will I submit the final secret proposal in the required format. This ensures the proposals are complementary and mutually accepted before the game ends.} \\
            }
        }
    }
     \\ \\

\end{supertabular}
}

\end{document}
